% Generated draft: requires mathematical and source review.
\documentclass[11pt]{article}
\usepackage[margin=25mm]{geometry}
\usepackage{fontspec}
\setmainfont{Latin Modern Roman}
\usepackage{amsmath,amssymb,mathtools,unicode-math}
\setmathfont{Latin Modern Math}
\usepackage{enumitem,needspace,xurl,hyperref}
\setlength{\parindent}{0pt}
\setlength{\parskip}{5pt}
\setlength{\emergencystretch}{3em}
\newcommand{\N}{\mathbb{N}}
\newcommand{\Z}{\mathbb{Z}}
\newcommand{\Q}{\mathbb{Q}}
\newcommand{\R}{\mathbb{R}}
\newcommand{\C}{\mathbb{C}}
\newcommand{\F}{\mathbb{F}}
\newcommand{\A}{\mathbb{A}}
\newcommand{\PP}{\mathbb{P}}
\newcommand{\E}{\mathbb{E}}
\newcommand{\eps}{\varepsilon}
\newcommand{\dd}{\,\mathrm{d}}
\begin{document}
\section*{UnsolvedMath ID: 30002155}
\textbf{Source problem code:} catalogue entry
\subsection{Definitions and mathematical statement}
The supplied record concerns the Breuil--M\'ezard multiplicity formula. Its intended setting is local $p$-adic Galois deformation theory: $p$ is a prime, $K$ is a finite extension of $\Q_p$, $G_K=\operatorname{Gal}(\overline K/K)$, and $\bar\rho:G_K\to\operatorname{GL}_2(\mathbb F)$ is a continuous residual representation over a finite field $\mathbb F$ of characteristic $p$. These are the conventional interpretations of the notation; the record does not explicitly specify all field and continuity hypotheses or any restrictions on $p$.

Let $\lambda$ denote a regular $p$-adic Hodge type and $\tau$ an inertial type. In the usual terminology, regularity means that the two labeled Hodge--Tate weights at each embedding are distinct; the record does not specify the labeling or sign convention. Write $\mathcal O_E$ for the ring of integers of a coefficient field $E$ and $\pi_E$ for a uniformizer. Let $R_{\lambda,\tau}$ denote the deformation ring intended by the source. For a nonzero Noetherian local ring $B$, write $e(B)$ for its Hilbert--Samuel multiplicity with respect to its maximal ideal. The precise deformation problem defining $R_{\lambda,\tau}$, and the convention for a zero deformation ring, are not supplied.

Let $A$ denote the source's unspecified index set for $a$. Write $n_a(\lambda,\tau)$ for the coefficients denoted $n(a)$ in the record, making explicit their intended dependence on the types. In the usual Breuil--M\'ezard formulation, these are multiplicities of residual weights in a representation associated with the types; the record does not identify that representation or its normalization.

For each fixed $\bar\rho$, the supplied question is whether there exists a family
\[
  \bigl(\mu_a(\bar\rho)\bigr)_{a\in A}\in\Z_{\geq 0}^{A}
\]
such that, for every regular Hodge type $\lambda$ and every inertial type $\tau$ in the intended deformation problem,
\[
 e\!\left(R_{\lambda,\tau}/\pi_E R_{\lambda,\tau}\right)
   =\sum_{a\in A}n_a(\lambda,\tau)\,\mu_a(\bar\rho).
\]
The same family must work simultaneously for all such pairs $(\lambda,\tau)$. The index set, coefficients, admissible types, and deformation rings must be specified before this identity becomes a fully determined mathematical conjecture; the record also does not explicitly state the finiteness conditions making the sum meaningful. In particular, it does not resolve whether the rings parametrize potentially semistable or potentially crystalline lifts, whether they are framed, or whether a determinant is fixed.
\subsection{Short English statement}
For a fixed residual two-dimensional Galois representation, ask whether every indicated deformation-ring multiplicity is a weighted sum of nonnegative integers depending only on the residual representation and the weight index. The coefficients may depend on the Hodge and inertial types. Several definitions needed to make this precise are missing from the supplied record. Its supplied status is partially solved, according to dated curation requiring independent verification.
\subsection{Sources}
\begin{itemize}
\item Supplied local UnsolvedMath registry record, permanent ID $30002155$, title ``Breuil--M\'ezard Multiplicity Formula,'' source problem code \texttt{catalogue entry}, in the Oberwolfach Reports Open Problems collection. No individual proposer is identified; the supplied proposed year is 2012.
\item Supplied source citation: \emph{Arithmetic Geometry (2013)}, Oberwolfach Reports, DOI 10.4171/owr/2012/38, pp.\ 2335--2388. \url{https://doi.org/10.4171/owr/2012/38}. The registry labels the statement \texttt{corrected\_verified} and reports that the official report gives the residual representation and quotient-ring multiplicity displayed here. This is supplied verification metadata, not an independent verification of the report.
\item Gee--Kisin, ``potentially Barsotti--Tate Breuil--Mezard'' (supplied reference label). \url{https://arxiv.org/abs/1208.3179}.
\item Caraiani--Emerton--Gee--Savitt, ``geometric Breuil--Mezard'' (supplied reference label). \url{https://msp.org/ant/2025/19-2/ant-v19-n2-p03-p.pdf}.
\item Supplied literature assessment, checked on exactly \texttt{2026-08-21}: the status is labeled partially solved; the full formula for arbitrary regular Hodge and inertial types over an arbitrary finite extension is described as not settled as one theorem, while potentially Barsotti--Tate and other geometric or special cases are described as proved. This is supplied dated curation derived from the OWR record, not independently verified literature status or confirmation of the cited claims.
\end{itemize}
\end{document}
